\section{Introduction}

This capstone project represents the culmination of several years of undergraduate engineering development, research, and high-power rocketry work conducted at the University of Hartford. The project originated through the Advanced Undergraduate Rocketry Operations and Research Association (AURORA), which was established to develop the technical capability, infrastructure, and flight experience required to design and operate large student-built research rockets.

Early AURORA efforts focused on developing vehicles for collegiate high-power rocketry competition. As project requirements, funding, and research priorities evolved, the effort was redirected toward the development of a reusable sounding-rocket platform. Hardware, flight experience, and engineering methods developed through earlier projects were retained and now form the foundation of the current capstone. The resulting platform is intended to support experimental avionics, telemetry, radio-frequency communication, recovery systems, and future student research.

The primary vehicle is the AURORA Rocket and Experimental Systems (\ares{}) platform, a 10.5-foot-tall, 6-inch-diameter high-power research rocket. \ares{} integrates mechanical design, aerodynamic and structural analysis, propulsion, recovery, embedded electronics, telemetry, RF systems, manufacturing, and flight operations within a single multidisciplinary vehicle. The rocket is being developed as a modular platform so that individual experimental and flight-critical systems can be independently developed, tested, and upgraded.

Three primary student-developed systems are incorporated into \ares{}. The Data Emitting \& Insight Module Orientation System (\deimos{}) is located within the nose cone and contains the sixth generation of the AURORA Insight telemetry platform. The current revision transitions previous development-board-based designs toward a fully embedded PCB with emphasis on reliability, compactness, higher data rates, extended communication range, and RF integration. The P.H.O.B.O.S. (\phobos{}) payload investigates electronically steered directional RF communication using phased-array antenna techniques, in which controlled phase relationships between antenna elements are used to alter the resulting radiation pattern \cite{balanis}. Flight-critical recovery operations are managed by the Compact Onboard Management for Ejection Timing (\comet{}) flight computer, a custom dual-deployment avionics system that uses barometric and inertial measurements to detect flight events and control the rocket's recovery system.

The project is supported through a combination of undergraduate research funding and prospective industry sponsorship. Development of the \phobos{} payload is funded through a NASA Connecticut Space Grant Consortium Undergraduate Research Grant awarded to Kyle Eckert during the Spring 2026 award cycle \cite{ctsgc2026}. The research objective has been retained as the overall vehicle architecture has evolved, allowing the RF experiment to be integrated into the current \ares{} platform. The \ares{} vehicle and \deimos{} avionics subsystem are also currently under consideration for sponsorship by Ensign-Bickford Aerospace \& Defense (EBAD), with prospective support directed toward vehicle materials, propulsion, electronics, manufacturing, integration, and testing.

The primary objective of the capstone is the design, manufacture, integration, verification, and flight testing of a large high-power research rocket incorporating both flight-critical and experimental student-developed systems. A successful mission will demonstrate that the independently developed mechanical, avionics, recovery, telemetry, and RF subsystems can operate together as a unified flight vehicle. Beyond the initial mission, \ares{} is intended to establish a reusable research platform and technical foundation for continued undergraduate aerospace development at the University of Hartford.